\documentclass[aspectratio=169,11pt]{beamer}

% Packages
\usepackage[utf8]{inputenc}
\usepackage[T1]{fontenc}
\usepackage{lmodern}
\usepackage{microtype}
\usepackage{amsmath,amssymb}
\usepackage{booktabs}
\usepackage{graphicx}
\usepackage{tikz}
\usepackage{pgfplots}
\usepackage{xcolor}
\usepackage{ragged2e}
\usepackage{colortbl}
\usepackage{array}
\usepackage{hyperref}
\usepackage{adjustbox}

\usetikzlibrary{shapes.geometric, arrows.meta, positioning, calc, backgrounds, decorations.pathreplacing, shadows, fadings}
\pgfplotsset{compat=1.18}

% --- Color Palette: Warm Professional ---
\definecolor{DeepNavy}{HTML}{2E4057}
\definecolor{Teal}{HTML}{048A81}
\definecolor{WarmOrange}{HTML}{E85D04}
\definecolor{SoftPurple}{HTML}{9D4EDD}
\definecolor{WarmGray}{HTML}{6C757D}
\definecolor{LightGray}{HTML}{E9ECEF}
\definecolor{Cream}{HTML}{FBF8F1}
\definecolor{DeepRed}{HTML}{D62828}
\definecolor{Gold}{HTML}{D4A03A}
\definecolor{SoftWhite}{HTML}{FAFAFA}

% Beamer colors
\setbeamercolor{normal text}{fg=DeepNavy, bg=SoftWhite}
\setbeamercolor{structure}{fg=DeepNavy}
\setbeamercolor{alerted text}{fg=DeepRed}
\setbeamercolor{example text}{fg=Teal}
\setbeamercolor{frametitle}{fg=DeepNavy, bg=SoftWhite}
\setbeamercolor{title}{fg=DeepNavy}
\setbeamercolor{subtitle}{fg=WarmGray}
\setbeamercolor{block title}{bg=DeepNavy, fg=SoftWhite}
\setbeamercolor{block body}{bg=LightGray, fg=DeepNavy}
\setbeamercolor{itemize item}{fg=WarmOrange}
\setbeamercolor{itemize subitem}{fg=Gold}

% Fonts
\usefonttheme{professionalfonts}
\setbeamerfont{title}{size=\huge, series=\bfseries}
\setbeamerfont{frametitle}{size=\Large, series=\bfseries}

% Strip chrome
\setbeamertemplate{navigation symbols}{}
\setbeamertemplate{headline}{}

% Minimal footline
\setbeamertemplate{footline}{%
    \hfill
    \begin{beamercolorbox}[wd=3cm,ht=2.5ex,dp=1ex,right,rightskip=0.6cm]{page number}
        \usebeamercolor[fg]{WarmGray}\scriptsize\insertframenumber
    \end{beamercolorbox}
    \vspace{0.3cm}
}

% Bullets
\setbeamertemplate{itemize item}{\footnotesize\raisebox{0.3ex}{\tikz\fill[WarmOrange] (0,0) circle (0.45ex);}}
\setbeamertemplate{itemize subitem}{\footnotesize\raisebox{0.3ex}{\tikz\fill[Gold] (0,0) circle (0.35ex);}}

% Custom commands
\newcommand{\transitionslide}[2]{%
    {\setbeamercolor{normal text}{fg=SoftWhite,bg=DeepNavy}
    \begin{frame}[plain]\vfill\begin{center}
        {\huge\bfseries\textcolor{Gold}{#1}}\\[0.6cm]
        {\large\textcolor{LightGray}{#2}}
    \end{center}\vfill\end{frame}}
}

\begin{document}

\begin{frame}
    \title{Appendix}
    \subtitle{Optimal Weights ($\Gamma$) and Factor Estimation}
    \author{Roy Kisluk}
    \date{}
    \titlepage
\end{frame}

\begin{frame}
    \frametitle{The Need for $\Gamma$: Leaky Instruments}
    
    \vspace{0.5cm}
    {\Large
        If large firms respond differently to macro trends, the simple GIV fails.
    }
    
    \vspace{1cm}
    \begin{itemize}
        \item \textbf{The Simple Model:} Assumes all firms have the same sensitivity ($\lambda_i = 1$) to macroeconomic shocks ($\eta_t$).
        \item \textbf{The Reality:} Superstar firms (e.g., Apple, Delta) often have very different macro-sensitivities than the average firm.
        \item \textbf{The Problem:} If big firms are more sensitive to $\eta_t$, then the size-weighted outcome ($y_{St}$) becomes highly correlated with the macro shock.
        \item \textbf{The Result:} The simple GIV instrument ($z_t = y_{St} - y_{Et}$) becomes "leaky"—it inadvertently instruments with aggregate shocks rather than just idiosyncratic ones.
    \end{itemize}
\end{frame}

\begin{frame}[shrink=20]
    \frametitle{Deriving $\Gamma$: Orthogonal Projection}
    
    \vspace{0.5cm}
    {\Large
        $\Gamma$ is the "pure granular weight", orthogonal to macro-sensitivities.
    }
    
    \vspace{0.5cm}
    We want weights $\Gamma_i$ that maximize power (by staying close to firm sizes $S_i$) while ensuring orthogonality to the common factors ($\lambda_i$) and a constant (to sum to zero).
    
    \vspace{0.5cm}
    \textbf{Optimization Problem:}
    \begin{equation*}
        \min_{\Gamma} \sum_{i=1}^N (\Gamma_i - S_i)^2 \quad \text{subject to} \quad \sum \Gamma_i = 0, \quad \sum \Gamma_i \lambda_i = 0
    \end{equation*}
    
    \vspace{0.5cm}
    \textbf{Solution (Linear Projection):}
    \begin{equation*}
        \Gamma = S - \lambda (\lambda' \lambda)^{-1} \lambda' S
    \end{equation*}
    \vspace{0.2cm}
    $\Gamma$ is simply the residual of regressing firm sizes ($S$) onto the factor loadings ($\lambda$). It removes the macro-contaminated portion of size.
\end{frame}

\begin{frame}[shrink=5]
    \frametitle{Estimating Latent Factors with PCA}
    
    {\large
        What if the true factor loadings ($\lambda_i$) are unobserved?
    }
    
    \vspace{0.1cm}
    In practice, we often do not know exactly how each firm responds to every macro trend. We must estimate the common factors ($\eta_t$) and their loadings ($\lambda_i$) from the data itself.
    
    \begin{itemize}
        \item \textbf{Principal Component Analysis (PCA):} We apply PCA to the panel data of firm outcomes ($y_{it}$) to extract the largest sources of common variance.
        \item \textbf{The First Principal Components:} The first $K$ principal components represent the dominant macroeconomic factors ($\eta_t$).
        \item \textbf{Factor Loadings:} The firm-specific sensitivities to these components become our estimated $\hat{\lambda}_i$.
        \item \textbf{Calculating $\Gamma$:} We then project firm size $S$ orthogonal to $\hat{\lambda}$ to find the optimal, bias-free GIV weights ($\Gamma$).
    \end{itemize}
\end{frame}

\end{document}